\documentclass[11pt,a4paper]{article}
%-------------------------
% Basic packages
%-------------------------
\usepackage[utf8]{inputenc}
\usepackage[T1]{fontenc}
\usepackage{lmodern}
\usepackage{microtype}
\usepackage{geometry}
\geometry{margin=1in}

% Mathematics
\usepackage{amsmath,amssymb,amsthm,mathtools,bm}
\usepackage{physics}    % convenient bra-ket, derivatives, etc.
\usepackage{siunitx}    % units
\usepackage{braket}     % alternate bra/ket macros
\usepackage{mhchem}     % chemistry notation if needed

% Figures & graphics
\usepackage{graphicx}
\usepackage{tikz}
\usepackage{pgfplots}
\pgfplotsset{compat=1.18}

% References, hyperlinks, and clever names
\usepackage[colorlinks=true,linkcolor=blue,citecolor=blue,urlcolor=blue]{hyperref}
\usepackage[nameinlink]{cleveref}

% Lists and layout
\usepackage{enumitem}
\setlist{itemsep=2pt,parsep=2pt}

% Theorem environments
\theoremstyle{plain}
\newtheorem{theorem}{Theorem}[section]
\newtheorem{lemma}[theorem]{Lemma}
\newtheorem{proposition}[theorem]{Proposition}
\newtheorem{corollary}[theorem]{Corollary}

\theoremstyle{definition}
\newtheorem{definition}[theorem]{Definition}
\newtheorem{example}[theorem]{Example}

\theoremstyle{remark}
\newtheorem{remark}[theorem]{Remark}

%-------------------------
% Useful macros
%-------------------------
% \newcommand{\dd}{\mathrm{d}}                      % differential
\newcommand{\ii}{\mathrm{i}}                      % imaginary unit
\newcommand{\eexp}{\mathrm{e}}                    % exponential base
\DeclareMathOperator{\sgn}{sgn}

% Shortcuts for common physics notation (overrides safe if physics package present)
\renewcommand{\vec}[1]{\boldsymbol{#1}}

%-------------------------
% Bibliography (biber/biblatex) - optional
%-------------------------
% \usepackage[backend=biber,style=phys]{biblatex}
% \addbibresource{references.bib}

%-------------------------
% Document metadata
%-------------------------
\title{Execise 6}
\author{Zhuge X.K.}
\date{\today}

%-------------------------
% Document
%-------------------------
\begin{document}

\maketitle
\section*{Exercise 2.7(E.2.5, page 115)}
Coherent inelastic transport. In this book we will generally restrict ourselves to steady-state transport in the persence of a d.c. bias. However, it is interesting to note that if an alternating field is present within the conductor (but not inside the contacts) then we can define scattering states just like those in Eq.2.6.3 but with different energies all coupled together. A scattering state $(q,E)$ now consists of an incident wave with energy $E$ in lead $q$, together with scattered waves with energy ($E_n = E \pm n\hbar \omega$, $n$ being an integer) in every lead $p$ (Eq.2.6.3)
\begin{equation*}
    \Psi_p(q) = \delta_{pq} e^{\ii k x_q}\exp(-\ii E t/\hbar) + \sum_{p,n} s'_{pq}(E_n,E) \exp(-\ii k_n x_p) \exp(-\ii E_n t/\hbar)
\end{equation*}
Proceeding as before derive the following expression for the d.c. current:
\begin{equation*}
    I_p = \frac{2e}{h} \int \left[ f_p(E) - \sum_{q,n}T_{pq}(E_n, E) f_q(E) \right] \dd E
\end{equation*}
The point is that there is no reason to include any exclusion principle-related factors in the expression for the current even if transport is inelastic, as long as it is coherent. But if there are phase-breaking processes within the conductor then we cannot define coherent scattering states that extend from one lead to another and the scattering state argument cannot be used.
\section*{Exercise 2.8(E.2.7, page 116)}
Starting from the current expression
\begin{equation*}
    I_p = -\frac{2e}{h} \int \overline{T}(E) \left[ f_1(E) - f_2(E) \right] \dd E
\end{equation*}
derive the result stated in Eqs.2.5.4-2.5.6.
\end{document}